\chapter{Anomaly Detection in Industrial Control Systems}
\label{chap:anomaly}

\section{Scope, Research Question, and Contributions}
\label{sec:anomaly:scope}
\label{sec:anomaly:1}

This chapter answers RQ3, which asks how a modular learning framework can analyse network and
process features for anomaly detection in industrial control systems
(Section~\ref{sec:intro:rqs}), and addresses gap~\gapref{3}. The framework, \nadics{}, supports
researcher-guided model development through configurable feature preparation and comparison of
interchangeable learning components. In the German Industrie~4.0
security project IUNO (Section~\ref{sec:anomaly:nadics}), its network application
uses a protocol-specific extractor for S7comm and PROFINET. A GAN-based extension analyses
multivariate SWaT process windows with elements of MAD-GAN and TadGAN.

The two industrial applications use different datasets and decision units. The network
evaluation classifies packet-derived connection records, while the SWaT study derives timestamp
decisions from process windows. They share a learning interface, but their results are assessed
under separate experimental protocols.

Released code and IUNO deliverables document the original framework \cite{iuno_ap4,iuno_ap7}.
The SWaT study is recorded in an unpublished internal report written by the candidate. Because
that report is a research record, it is not a bibliography item. The candidate designed the
method, implemented the experimental pipeline, ran and analysed the experiments, and wrote the
report. Section~\ref{sec:anomaly:gan} presents its empirical evaluation.

Section~\ref{sec:anomaly:service-graphs} retains a short exploratory IoT study as an example of
relation-based monitoring \cite{aubet2018graph}. It was separate from \nadics{} and provides
supporting context; the answer to RQ3 rests on the industrial framework and its evaluations.

\section{Observation Models and Datasets}
\label{sec:anomaly:observation}

\nadics{} (Network Anomaly Detection for Industrial Control Systems) is the detection framework
described in Section~\ref{sec:anomaly:nadics} \cite{iuno_ap4}. It receives tabular records. Its
network path takes connection features derived from the industrial protocols S7comm and PROFINET,
or from conventional network traffic. The SWaT path consumes the sensor and actuator values of
a six-stage water-treatment process, sampled once per second
\cite{goh2016swatdataset,mathur2016swat}.

\begin{table}[tbp]
  \centering
  \caption{Observation units in the industrial evaluations and the exploratory IoT study.
  Labels, sampling processes, and validation protocols differ between the studies.}
  \label{tab:anomaly:observations}
  \small
  \begin{tabularx}{\textwidth}{@{}>{\raggedright\arraybackslash}p{2.5cm}
    *{3}{>{\raggedright\arraybackslash}X}@{}}
    \toprule
    Study & Unit observed & Representation & Decision \\
    \midrule
    \nadics{} in IUNO &
      packet-derived connection record &
      protocol and flow features &
      binary or attack-class label \\
    SWaT GAN study &
      sensor and actuator sample &
      multivariate time window &
      anomalous or normal timestamp \\
    IoT service graph (exploratory) &
      inter-service message &
      vertices and whitelisted relations &
      known or unknown relation \\
    \bottomrule
  \end{tabularx}
\end{table}

The industrial evaluations use different decision rules. The S7-300 evaluation
(Section~\ref{sec:anomaly:iuno-corpus}) uses labelled classification. The GAN study
(Section~\ref{sec:anomaly:gan}) fits mixed normal and attack windows, combines reconstruction
and critic scores, and applies an adaptive threshold to assign a binary prediction to each
timestamp. An attack present during fitting may enter its learned reference distribution.
Table~\ref{tab:anomaly:observations} also includes the service-relation unit of the exploratory
study below.

\section{Exploratory Study: Communication Graphs for IoT Services}
\label{sec:anomaly:service-graphs}

The collaborative study presented at PAM 2018 examined communication whitelisting for IoT
microservices \cite{aubet2018graph}. Its monitor operates at the Knowledge Agent, the entry
point of the Virtual State Layer middleware through which services exchange data
\cite{pahl2016distributed}. Each computing node stores a graph whose vertices are service
addresses and whose edges record observed communication partners. This observation model
captures service relations at the middleware level, in contrast to the packet-derived records
used in the industrial evaluation.

A learning phase begins when a service starts and treats the relations observed during that
period as normal. A later relation absent from the graph raises a notification. The site owner
can approve it, after which the monitor adds it to the allowed relations. This update procedure
permits new legitimate partners during operation, subject to the assumption that the initial
learning period represents benign behaviour. The prototype uses concrete local addresses;
site-independent service models remained a proposed extension \cite{aubet2018graph}.

The source reports a data-structure memory estimate of 33.7~KB for twenty services with ten
relations each. Measurements with ten to forty services give an added latency of about 6~ms
over an inter-service round trip of 24~ms, although the repetition count and the summary
statistic behind these values are unspecified. A short graph-update trace records service
start-up and the addition of a service, followed by a notification for a previously unseen
connection \cite{aubet2018graph}.

The evidence documents prototype operation and resource costs. Detection quality was not
measured against a labelled attack set, so the study provides neither attack recall nor a
false-positive rate. The decision rule identifies unfamiliar communication partners. A new
legitimate relation may therefore raise an alert, while malicious activity that preserves the
permitted relations remains outside the rule's scope.

\section{The NADICS Framework and IUNO Integration}
\label{sec:anomaly:nadics}
\label{sec:anomaly:5}

\nadics{}, Network Anomaly Detection for Industrial Control Systems, is a Python framework
built at the Chair of IT Security of the Technical University of Munich to detect attacks
against PLC and SCADA installations from network traffic. The candidate developed it to help
researchers construct a suitable detection pipeline for a new environment or dataset. Its
design goal was semi-automated model development: researchers define the prediction task and
candidate configurations, while the framework coordinates repeated training and comparison.
The intended selection process includes data preparation and hyperparameter tuning alongside
the choice of learning algorithm.

The framework is organised in three modules. Preprocessing prepares the feature table through
encoding and scaling, with optional resampling for class imbalance. Feature-importance analysis
guides the researcher's choice of input variables. The learning module holds the estimators behind
one interface, so that classical models and neural models can be exchanged by name without the
pipeline being rewritten: the released repository contains the scikit-learn estimators, the IUNO
deliverable additionally reports a deep neural network and a convolutional network
\cite{iuno_ap4}, and the generative adversarial detector of Section~\ref{sec:anomaly:gan} was
added later at the same interface for multivariate process data. Postprocessing produces
evaluation reports and supplies detections to site-specific alarm handling. The S7comm and
PROFINET extractor existed as a separate project
artefact, so the released repository takes a feature table as its input.

The IUNO deliverable records the need to produce labelled traffic from a real industrial
environment \cite{iuno_ap4}. For a new dataset, researchers must establish its feature meanings
and the available labels before choosing compatible methods. The layout and model registries
allow these choices to change without a separate training and reporting script for each
algorithm. The objective is to identify a configuration suited to that dataset and a stated
criterion, such as attack recall subject to a tolerable false-positive rate. The preferred
configuration is conditional on the candidates considered and the evaluation protocol.

\begin{figure}[p]
  \centering
  \includegraphics[width=\linewidth, trim=5 157 5 5, clip]%
    {nadics_arch_IUNO_overall_pipeline.pdf}
  \caption{Architecture of the \nadics{} framework. Malicious traffic is generated in a real
  industrial environment, dissected into protocol-characteristic features, and classified by the
  anomaly detection engine. The stage structure follows the one reported in the IUNO
  deliverable \cite{iuno_ap4}.}
  \label{fig:anomaly:nadics-arch}
\end{figure}

IUNO, the German national reference project for IT security in Industrie~4.0, ran from 2015 to
2018. Its consortium joined industrial companies such as Volkswagen with research institutes,
among them Fraunhofer AISEC, Fraunhofer SIT, DFKI and TUM \cite{iuno2017d42}. The partners
built security tools and tested them in four use cases \cite{iuno_ap7}. \nadics{} was developed
for the use case \emph{Visual Security Control Room}. This control room captured data streams
from production in real time and displayed the security state of the production network in one
view, so that operators could act on a deviation before the process came to a halt
\cite{iuno_ap4}.

Figure~\ref{fig:anomaly:nadics-arch} shows the pipeline as it was deployed in that use case. Two
data artefacts separate its three stages. The first stage produced the corpus. Attacks of five
classes (authentication bypass, CPU stop and start, reconnaissance, denial of service and brute
force) were executed in a real industrial environment against Siemens S7-300 controllers that
exchanged S7comm traffic \cite{kleinmann2014accurate} with a control server. The traces were
labelled both in binary form and by attack class \cite{iuno_ap4}. In the second stage, the
public analysers TShark, Bro-IDS and Argus were combined with dissectors written for S7comm and
PROFINET
\cite{akerberg2009exploring}, and the feature extractor reduced each network connection to
one labelled feature vector, written as CSV or JSON
\cite{iuno2017d42,iuno_ap4}. The third stage is the framework proper. It reads that dataset, fits
one of the registered estimators, applies the resamplers and the outlier tests, and reports which
features carried the decision. Detections reached the control room as JSON event messages over
the authenticated message bus of the IUNO aggregation system \cite{iuno_ap4}. The response to
an alert was left to the control room operators \cite{iuno_ap7}.

The modules map onto these stages at the dataset. Preprocessing spans the second stage and the
input side of the third, while learning and postprocessing both sit in the third stage. Corpus
generation lies outside the framework. A layout file specifies the input columns and target
labels, together with the feature subset and the string fields that require encoding. Model
registries map names to estimators. The original repository contains thirteen classifiers and
four novelty or outlier detectors, with fifteen resamplers and two boosting meta-estimators.
The registry named \texttt{clustering} supplies the novelty and outlier methods. Supervised
learners use \texttt{fit} and \texttt{predict}; the outlier routines also use score functions
or \texttt{fit\_predict}, according to the estimator interface \cite{pedregosa2011scikit}.

The archived runner accepts a researcher-specified list of classifiers and reports their
predictive scores together with training and prediction times. Researchers can request
feature-importance estimates from an Extra-Trees ensemble and revise the selected columns in
the layout. The \texttt{learning.py} module defines a \texttt{GridSearchCV} helper with
five-fold cross-validation over predefined parameter grids for decision trees, random forests
and $k$-nearest neighbours, with scoring left at the estimator default. These components
document the intended automation of experiment execution and bounded parameter search.
% Source inspection: NADICS_repo/nadics/{main,models,selector,learning,parser,report}.py.

Model selection for a new environment requires a development protocol separate from final
testing. Each validation split must fit data-dependent preprocessing and feature selection on
its training portion; resampling must leave validation observations unchanged. After selecting
a configuration, researchers assess it on held-out data and define how its predictions should
be converted into operational alerts. This workflow reuses the analysis components while
retaining domain-specific decisions about labels and acceptable alarm burden. The industrial
experiments below evaluate configured detectors; they do not measure an advantage of automated
pipeline selection over manual development.

\begin{sidewaysfigure}
  \centering
  \includegraphics[width=\textheight]{iuno-VW_testbed.pdf}
  \caption{The IUNO test environment at the Volkswagen plant in Wolfsburg. Four servers, three
  network segments and a Siemens plant network reproduce the production setting in which
  \nadics{} ran as the TUM detection tool on the virtualisation host. The network TAPs fitted at
  the relevant interfaces are omitted from the drawing \cite{iuno2017d42}.}
  \label{fig:anomaly:iuno-testbed}
\end{sidewaysfigure}

Figure~\ref{fig:anomaly:iuno-testbed} shows where the framework ran. The environment was built
at the Volkswagen plant in Wolfsburg from hardware of the kind used in regular production
\cite{iuno2017d42}. A Hirschmann industrial switch carries three segments, the intranet, the
demilitarised zone and the plant network, and manually configured Eagle routers separate them.
The first of four servers runs the WinCC and OPC~UA clients that originate test traffic, together
with a Nagios instance and the database of the aggregation system. The second is the
virtualisation host for the detection tools contributed by AISEC, DFKI, SIT and TUM, and the
third serves the visual control room to the intranet. The fourth combines the WinCC server with
the collector, which gathers traffic from the plant network and passes it, bundled and
pre-filtered, to the detection tools, so \nadics{} was one of four detectors reading the same
collector output. Siemens controllers occupy the plant network behind a Scalance switch. The
collector observed that network passively at the plant switch, never in line with the traffic,
which matches the detection-only design: an alert could reach the control room, but nothing in
the path could stop a packet.

The IUNO deliverables describe the tool and its environment but print no measured result for it
\cite{iuno_ap4}. The values below come from the experimental records of the project and cover
the corpus and its feature schema, the detection results on the held-out session, the alarm
burden over a benign monitoring period, and the throughput of the extractor on the prototype
collector.

\subsection{The S7-300 Corpus}
\label{sec:anomaly:iuno-corpus}

The corpus was recorded in 2018 in the environment of Figure~\ref{fig:anomaly:iuno-testbed},
where six Siemens S7-300 controllers exchanged routine polling and control traffic with
WinCC-based services. Complete selected channels were kept, together with ARP and DCP traffic and
the connection attempts belonging to those channels. No random packet sampling was applied.

Recording opened with a two-hour clean baseline, which was excluded from the corpus totals. That
interval fixed the communication relationships, the polling behaviour and the protocol activity
expected under the operating modes present at the time. Six capture sessions followed, with 24
hours of labelled traffic in total. The collector recorded 12,616,875 Ethernet frames in that
period, an average of 146.03 frames per second. The captured frames hold 2,317,074,962 bytes,
and the six capture files occupy about 2.519~GB in the classic PCAP format.

Twenty-five controlled attack episodes were executed during the sessions, five repetitions of
each of the five classes the deliverable names \cite{iuno_ap4}. Background communication
continued while they ran. Table~\ref{tab:anomaly:iuno-episodes} lists the duration of every
episode. The shortest lasted 3~min 07~s and the longest 9~min 55~s, and together they cover
2~h 51~min 05~s. Episodes also differed in the channels they touched and in how often a
connection was re-established. Labels were assigned
from the execution logs and from inspection of the packet traces, and only the communications an
episode actually affected received an attack label, so background traffic concurrent with an
attack keeps the normal label. Two attack types were never run against one channel at the same
time.

\begin{table}[tbp]
  \centering
  \caption{Composition of the S7-300 corpus and its partitions. Whole capture sessions were
  assigned to the partitions, four to training and one to each of the other two, and every
  attack episode lies entirely within one partition. The episode column counts executed attacks,
  not records.}
  \label{tab:anomaly:iuno-corpus}
  \small
  \begin{tabular}{@{}lrrrrr@{}}
    \toprule
    Class & Corpus & Train & Validation & Test & Episodes \\
    \midrule
    Normal                 & 113,751 & 76,556 & 18,587 & 18,608 &    \\
    Authentication bypass  &   1,058 &    441 &    269 &    348 &  5 \\
    CPU stop and start     &     593 &    359 &     90 &    144 &  5 \\
    Reconnaissance         &   1,803 &  1,070 &    366 &    367 &  5 \\
    Denial of service      &   1,230 &    765 &    278 &    187 &  5 \\
    Brute force            &     725 &    445 &     93 &    187 &  5 \\
    \midrule
    Total                  & 119,160 & 79,636 & 19,683 & 19,841 & 25 \\
    Capture time (h:min:s) & 24:00:00 & 16:01:57 & 3:56:27 & 4:01:36 & \\
    \bottomrule
  \end{tabular}
\end{table}

\begin{table}[tbp]
  \centering
  \caption{Duration of the 25 attack episodes in minutes and seconds. Each attack type has three
  training episodes, one validation episode and one test episode.}
  \label{tab:anomaly:iuno-episodes}
  \small
  \begin{tabular}{@{}llrrr@{}}
    \toprule
    Class & Training episodes & Validation & Test & All five \\
    \midrule
    Authentication bypass & 5:33, 4:01, 5:50 & 7:16 & 9:34 & 32:14 \\
    CPU stop and start    & 7:03, 8:23, 3:07 & 4:04 & 7:39 & 30:16 \\
    Reconnaissance        & 6:42, 9:40, 7:28 & 9:55 & 8:27 & 42:12 \\
    Denial of service     & 6:41, 9:39, 7:07 & 7:46 & 5:38 & 36:51 \\
    Brute force           & 8:05, 3:16, 4:06 & 4:15 & 9:50 & 29:32 \\
    \midrule
    Total                 &                  &      &      & 171:05 \\
    \bottomrule
  \end{tabular}
\end{table}

Extraction produced 119,160 feature records. Of these, 113,751 are normal and 5,409 carry one of
the five attack labels, an attack share of 4.54 per cent. Table~\ref{tab:anomaly:iuno-corpus}
gives the composition. The training sessions hold 79,636 records, the validation session 19,683
and the test session 19,841. Each attack type contributes three training episodes and one
episode to each of the remaining partitions, and because episodes differ in duration and in the
channels they touch, the attack records of a class are not divided between the partitions in a
fixed proportion. An episode, its recovery interval and the windows that depend on it stay inside
a single partition, and no rolling window crosses a partition boundary. The attack counts must
not be read as a count of attacks. A single episode of a few minutes generates many records,
and the corpus averages about 216 affected records per episode.

The extraction schema has 63 fields in five groups: 18 flow fields, 13 timing fields, 18
protocol fields, nine relational fields and five process fields. Forty-six are numerical, eleven
categorical and six binary, and encoding of the categorical fields expands the input
dimensionality of the model. The flow group follows the output of Argus and Bro-IDS, which the
extractor wrapped. Protocol fields cover S7 message and function codes,
addressed memory areas, address entropy and error codes, together with COTP and OPC~UA fields.
Relational fields describe the communication partners, the age of an edge, and the number of
distinct destinations and ports observed in the preceding minute. The process group reads one
preconfigured anchor tag per channel, whose identity and units are held as metadata, so a value
delta is never computed across unrelated tags. Unobservable protocol fields carry an explicit
missing state, and an absent OPC~UA security policy is not encoded as the policy \emph{None}.

\begin{figure}[tbp]
  \centering
  \includegraphics[width=\linewidth]{nadics-feature-importance.pdf}
  \caption{Random-forest importance of the 63 extraction fields. (a) The fifteen most
  important fields, coloured by group. (b) Each group's share of the total importance
  against its share of the field count. (c) Cumulative share over the ranked fields.}
  \label{fig:anomaly:iuno-importance}
\end{figure}

Figure~\ref{fig:anomaly:iuno-importance} shows the importance of each field in the random forest
of Section~\ref{sec:anomaly:iuno-results}. The forest was fitted to the training partition; the
impurity reductions from splits on each field were summed across its trees, normalised over all
63 fields, and ranked. Panel~(a) shows the fifteen highest, and
Appendix~\ref{app:nadics-importance} lists all 63. \texttt{s7\_function} leads at 11.9 per cent,
because a stop or a start is an ordinary S7 operation whose function code is visible in the
request, and \texttt{role\_violation} follows at 8.0 per cent, because a write from a host whose
role does not write is anomalous even over an established connection. Both signals have limits:
legitimate maintenance issues the same function codes, and a compromised engineering station
writes from a role that is permitted to write.

Panel~(b) sets each group's share of the total importance against its share of the fields.
Protocol semantics holds the largest share at 36.4 per cent, although it also holds the most
fields. The comparison that carries information is the share per field: the nine relational
fields average 2.38 per cent each, against 0.74 for the eighteen flow fields and 0.60 for the five
process fields. Only the process fields receive less importance per field than the flow
statistics. That is consistent with the position the schema takes: an
industrial detector built on Argus and Bro flow features alone discards most of what separates an
S7 attack from S7 polling. Panel~(c) gives the cumulative share. Ten fields reach 54.5 per cent
of the total importance, twenty-four reach 80 per cent and thirty-six reach 90 per cent, so the
tail of the schema is thin and a reduced schema is plausible. \nadics{} exposes an importance
ranking for exactly this purpose \cite{iuno_ap4}.

One record covers a non-overlapping ten-second window of an activity-bearing connection or
channel, and a shorter final record is written when a session closes. PROFINET RT traffic is
keyed at layer 2, since it carries no transport-layer connection. The windowing refines the
per-connection unit of Figure~\ref{fig:anomaly:nadics-arch} so that a change during a long-lived
session can be scored while the session is still open. Channels that fall silent emit no record.

Three definitions were adjusted so that no record depends on information from after its own
window closes. The duration field measures the exported segment and not the eventual lifetime of
an open connection, while a separate field retains the time elapsed since the underlying
connection opened. The ratio of unanswered requests counts only those requests whose one-second
response deadline had already expired at window closure, and pending requests are excluded from
both the numerator and the denominator. Trailing source and timing state is held over 60
seconds, and the period features are computed from past observations. Encoders, imputation and
scaling were fitted on the training partition alone, and the classifier together with its
preprocessing was then held fixed for validation and testing.

\subsection{Detection Results}
\label{sec:anomaly:iuno-results}

The classifier is a random forest of 300 trees with a maximum depth of 16 and a minimum leaf
size of five. Class weighting was applied per bootstrap sample, and the training data were not
resampled. Two decisions were read from the same fitted model: a binary verdict of normal
against attack, and an assignment to one of six classes, the normal class and the five attack
types. Every figure below comes from the held-out test session, which holds 18,608 normal and
1,233 attack records.

\begin{table}[tbp]
  \centering
  \caption{Six-class confusion matrix on the 19,841-record test session, with rows giving the
  true class. Predicted classes are abbreviated: N normal, AB authentication bypass, CS CPU stop
  and start, RC reconnaissance, DS denial of service, BF brute force. The last three columns
  give precision, recall and $F_1$ for the class of that row.}
  \label{tab:anomaly:iuno-confusion}
  \small
  \setlength{\tabcolsep}{4.5pt}
  \begin{tabular}{@{}lrrrrrrrrr@{}}
    \toprule
    & \multicolumn{6}{c}{Predicted class} & \multicolumn{3}{c}{Score (\%)} \\
    \cmidrule(lr){2-7}\cmidrule(l){8-10}
    True class & N & AB & CS & RC & DS & BF & Prec. & Rec. & $F_1$ \\
    \midrule
    Normal                & 18,518 &  21 &   5 &  31 &  18 &  15 & 99.74 & 99.52 & 99.63 \\
    Authentication bypass &     22 & 292 &   8 &  12 &   1 &  13 & 88.75 & 83.91 & 86.26 \\
    CPU stop and start    &      3 &   1 & 134 &   2 &   3 &   1 & 90.54 & 93.06 & 91.78 \\
    Reconnaissance        &      9 &   7 &   0 & 337 &   6 &   8 & 84.25 & 91.83 & 87.87 \\
    Denial of service     &      5 &   0 &   0 &   4 & 178 &   0 & 85.99 & 95.19 & 90.36 \\
    Brute force           &      9 &   8 &   1 &  14 &   1 & 154 & 80.63 & 82.35 & 81.48 \\
    \bottomrule
  \end{tabular}
\end{table}

Table~\ref{tab:anomaly:iuno-confusion} reports the six-class confusion matrix together with the
per-class scores. Collapsing the same predictions into the binary decision gives 18,518 true
negatives, 90 false positives, 48 false negatives and 1,185 true positives. Binary accuracy is
99.30 per cent, attack precision 92.94 per cent and attack recall 96.11 per cent, for an attack
$F_1$ of 94.50 per cent. The false-positive rate over normal records is 0.484 per cent and the
balanced accuracy 97.81 per cent. On the six-class task, accuracy reaches 98.85 per cent, the
macro-$F_1$ 89.56 per cent and the weighted $F_1$ 98.86 per cent.

Accuracy is the least informative of these numbers. Normal records make up 93.79 per cent of the
test partition, so a detector that never fired would already score 93.79 per cent, and the gap
between the weighted and the macro $F_1$, 98.86 against 89.56 per cent, measures how much of the
six-class figure the normal class carries. Per-class recall is highest for denial of service at
95.19 per cent and lowest for brute force at 82.35 per cent, with authentication bypass close
behind at 83.91 per cent. Most errors of these two classes fall on other attack classes.
Authentication bypass loses 22 of its 348 records to the normal class and 34 to the other attack
types, while brute force loses nine of its 187 records to the normal class and 24 to the other
attack types, 14 of them to reconnaissance. Both attacks reuse the protocol operations of
legitimate engineering access, so a window holding only a few such operations offers little that
separates them from ordinary traffic or from each other.

All of these scores are computed over records. They show that the extracted features separate
the five executed attacks from the surrounding traffic of this plant under a session-disjoint
split.

\subsection{Benign Operation and Alarm Burden}
\label{sec:anomaly:iuno-field}

A second exercise measured what the detector does when nothing malicious happens. It ran in the
Volkswagen environment with four active PLCs and ordinary supervisory traffic. A separate
two-hour baseline preceded 24 hours of benign monitoring. Over that day the collector recorded
8,472,913 frames, an average of 98.07 frames per second. The captured frames hold 2,186,039,471
bytes, about 2.322~GB in the classic PCAP format, and the extractor emitted 83,847 records.
Engineering work and planned device restarts took place during the period and count as benign.

The detector classified 83,740 records as normal and 107 as attacks, all of them false, a
record-level false-positive rate of 0.128 per cent. Positives were merged into alarm episodes
before presentation to the control room: positives that share a source, a destination and a
predicted attack class belong to one episode, and an episode closes once 60 seconds pass without
a further positive. One notification opens an episode, and later positives update it. Nine alarm
episodes were raised over the day, 0.375 per hour on average, and each was assigned a benign
cause. Engineering workstation activity accounts for 43 positive records in three episodes and
OPC~UA reconnections for 28 records in two. Brief communication interruptions produced 23
records in two episodes, and planned device restarts 13 records in two.

No attack was executed during this period, so every alarm it raised was false. Its 0.128 per
cent record rate lies below the 0.484 per cent of the test session, whose normal records also
include the recovery intervals of executed attacks. The recovered causes share a structure that
matters for deployment: each is a legitimate departure from routine polling, so the alarm burden
of this design follows the rate at which such events occur in the plant.

\subsection{Extractor Throughput on Constrained Hardware}
\label{sec:anomaly:iuno-pi}

The extractor was also exercised on a Raspberry Pi 3 Model B. The deliverable of June 2017
names Raspberry Pi boards as the prototype platform \cite{iuno2017d42}. The board carries four
cores at 1.2~GHz, 1~GB of RAM and 100~Mbit/s Ethernet. Only the extraction stage ran on it,
comprising selected-field capture and dissection of S7comm and PROFINET traffic, the bounded rolling
feature computations and emission of CSV or JSON. Model fitting and inference remained on the
server, and archival of full packet captures to the SD card was excluded. Each load replayed
traffic over 100 configured communication channels with varying channel activity for about ten
minutes.

\begin{table}[tbp]
  \centering
  \caption{Extraction workloads on the Raspberry Pi 3 prototype. Input frames are the frames
  presented to the extraction pipeline, and processed frames exclude those dropped before
  processing. Rates and CPU utilisation are averages over the replay, and CPU utilisation refers
  to the capacity of the whole device.}
  \label{tab:anomaly:iuno-pi}
  \small
  \begin{tabular}{@{}lrrr@{}}
    \toprule
    Quantity & Light load & Moderate load & Overload \\
    \midrule
    Replay duration (s)                 &  600.14 &    599.87 &    600.09 \\
    Input frames                        & 599,827 & 1,498,463 & 2,100,731 \\
    Dropped frames                      &       4 &       361 &   245,897 \\
    Processed frames                    & 599,823 & 1,498,102 & 1,854,834 \\
    Frame loss (\%)                     & 0.00067 &   0.02409 &  11.70531 \\
    Input rate (frames/s)               &  999.48 &  2,497.98 &  3,500.69 \\
    Processed rate (frames/s)           &  999.47 &  2,497.38 &  3,090.93 \\
    Feature records                     &   5,842 &     5,887 &     5,781 \\
    \midrule
    Mean CPU utilisation (\%)           &    22.6 &      43.1 &      50.4 \\
    Peak extractor memory (MiB)         &     302 &       381 &       436 \\
    95th-percentile emission delay (ms) &      96 &       234 &       817 \\
    \bottomrule
  \end{tabular}
\end{table}

Table~\ref{tab:anomaly:iuno-pi} reports the three loads. At the moderate load, 1,498,463 frames
were presented in 599.87 seconds and 361 were dropped, a loss of 0.024 per cent at a processed
rate of 2,497.38 frames per second. This input rate is about 17 times the average frame rate of
the corpus capture. The extractor used 43.1 per cent of the CPU capacity of the device on average
and at most 381~MiB of memory. The 95th percentile of the delay from the close of a window to the
emission of its record was 234~ms, a delay that excludes the ten seconds in which the window
accumulates. The light load lost four of its 599,827 frames at 22.6 per cent CPU utilisation.
The overload row does not describe a supported rate: 11.71 per cent of the input was dropped,
and the processed rate stayed at 3,090.93 frames per second against an input rate of 3,500.69.
Mean CPU utilisation reached 50.4 per cent there, and the 95th-percentile emission delay rose to
817~ms. The three runs emitted 5,842, 5,887 and 5,781 records, fewer than the 6,000 that 100
continuously active channels would give in ten-second windows over ten minutes, because channel
activity varied.

\section{A GAN-Based Multivariate Detector for SWaT}
\label{sec:anomaly:gan}

Industrial process recordings form multivariate time series: each timestamp associates sensor
measurements with actuator states, and successive observations describe the evolution of the
plant. An anomaly may appear as an unusual value or as a sequence that conflicts with the
expected relationships between variables. A measurement within its usual range can therefore
be anomalous in its temporal or process context. MAD-GAN motivates joint modelling of these
dependencies for industrial intrusion detection, while TadGAN addresses anomalous sequences
across several time-series domains \cite{li2019madgan,geiger2020tadgan}. Process monitoring can
expose the physical effects of an attack, although a deviation alone does not establish its
cause. Scarce labelled attack examples motivate methods that learn reference behaviour from
recorded operation.

A generative adversarial network (GAN) learns a data distribution through two neural networks
with opposing training objectives \cite{goodfellow2014generative}. The generator transforms
samples from a latent distribution into synthetic observations. The discriminator learns to
distinguish recorded observations from generated samples, and the generator is updated to make
that distinction harder. For time series, the generated objects are sequences of measurements.
Here, \emph{adversarial} describes the training relationship between the networks. The
discriminator's training labels identify recorded and generated data; they are not attack
annotations.

For anomaly detection, the learned distribution provides a reference against which new
observations can be assessed. MAD-GAN uses long short-term memory (LSTM) networks to model
multivariate windows and combines discrimination with reconstruction error
\cite{li2019madgan}. Reconstruction requires a latent representation from which the generator
can reproduce the observed sequence; MAD-GAN searches for this representation at inference.
TadGAN learns an encoder for this mapping and adds a cycle-consistency loss that penalises
differences between an input and its reconstruction. Its Wasserstein critics supply real-valued
scores for matching distributions during training \cite{geiger2020tadgan}. Reconstruction error
and discriminator or critic output can then supply evidence of deviation. Their usefulness
depends on the behaviour represented in training and on the subsequent scoring rule.

The present study combines multivariate process windows with the encoder and cycle-consistent
training of TadGAN. It examines how five input representations affect detection and false alarms
on a fixed SWaT evaluation cohort. The experiment concerns a detector developed with some
attack observations, as specified below; its learned reference therefore includes both normal
and attack behaviour. This distinction matters when comparing its results with methods trained
only on normal operation.

The detector operates on process measurements from the Secure Water Treatment (SWaT) testbed.
SWaT comprises six treatment stages, with programmable logic controllers coordinating sensors
and actuators and a supervisory network connecting the controllers to supervisory control and data acquisition (SCADA)
(Figure~\ref{fig:anomaly:swat-architecture}). The process dataset contains 51 sensor and
actuator variables recorded at a nominal frequency of one observation per second
\cite{goh2016swatdataset}. These variables describe the physical behaviour of the plant,
allowing the detector to consider changes in individual measurements together with relationships
among measurements across the treatment process.

\begin{figure}[htbp]
  \centering
  \includegraphics[width=\linewidth]{swat/thesis/swat-ics-architecture.pdf}
  \caption{SWaT process and control architecture. Network levels follow the SWaT convention.
  The field-network symbols indicate sensors and actuators within each stage; the detector
  consumes their recorded values. Adapted from the supplied diagram and the dataset
  description \cite{goh2016swatdataset}.}
  \label{fig:anomaly:swat-architecture}
\end{figure}

The dataset documentation describes 36 attacks distributed across four categories
(Figure~\ref{fig:anomaly:swat-taxonomy}): 26 single-stage, single-point attacks; four
single-stage, multiple-point attacks; two multiple-stage, single-point attacks; and four
multiple-stage, multiple-point attacks \cite{goh2016swatdataset}. These categories describe the
scope of the interventions. They do not imply that all attacks produce equally observable process
changes. An intervention affecting one sensor or actuator may also influence measurements
elsewhere in the plant, while an attack with little physical effect may remain difficult to
identify from process data alone. The aggregate calculations in this section concern anomalous
timestamps and do not establish the number of individual attacks detected.

\begin{figure}[htbp]
  \centering
  \includegraphics[width=\linewidth]{swat/thesis/swat-attack-taxonomy.pdf}
  \caption{SWaT attack categories and their counts in the dataset paper
  \cite{goh2016swatdataset}. Stage and point locations are schematic examples. The
  arrangement describes attack scope and carries no ranking of detection difficulty.}
  \label{fig:anomaly:swat-taxonomy}
\end{figure}

The source collection campaign comprises seven days of normal operation followed by four days
containing attacks \cite{goh2016swatdataset}. Preprocessing removes the first 21,600 observations
from the normal-operation trace, which cover the initial six-hour stabilisation period
(Figure~\ref{fig:anomaly:swat-campaign}). The empirical study uses all 449,919 timestamps of the
attack-period recording. Its first 56,240 timestamps are divided between the training and
validation partitions, which also receive separate portions of the normal-operation recording.
The later 393,679 attack-period timestamps form the test set. They comprise 51,453 anomalous and
342,226 normal timestamps, so anomalous observations account for 13.07\% of the test cohort.

\begin{figure}[htbp]
  \centering
  \includegraphics[width=\linewidth]{swat/thesis/swat-collection-campaign.pdf}
  \caption{SWaT collection campaign and the partition used in the GAN study. The campaign
  follows the dataset paper \cite{goh2016swatdataset}. After the initial normal-operation
  exclusion, the first 56,240 attack-period timestamps support training and validation; the
  later 393,679 timestamps form the test set.}
  \label{fig:anomaly:swat-campaign}
\end{figure}

All three partitions contain normal observations and labelled attack examples. Each attack
episode is confined to one partition, and no window crosses a partition boundary or a
discontinuity in the recordings. Partitioning precedes imputation, scaling, feature selection
and window construction. Earlier observations support model development, while the later
observations form the test cohort.

The GAN objective does not require class annotations during gradient-based fitting.
Nevertheless, the complete procedure includes supervised feature selection and model
selection, and training includes attack observations. The resulting evaluation therefore
concerns a detector developed with prior exposure to attacks. It does not establish detection
performance for previously unseen attack types or a wholly unsupervised development process.

\subsection{Model and Training}
\label{sec:anomaly:gan-method}

The model contains four jointly trained networks: an encoder, a generator, a data critic and a
latent critic (Figure~\ref{fig:anomaly:gan-pipeline}). For an input window $X$, the encoder $E$
maps the observed process sequence to a latent representation $E(X)$. The generator $G$ maps
this representation back to the observation space, producing the reconstruction
\[
  \widehat{X}=G(E(X)).
\]
Latent samples $z$ are drawn from a standard normal distribution. The data critic compares
observed windows with generated windows $G(z)$, encouraging the generator to reproduce the
distribution of the training sequences. The latent critic compares encoded observations $E(X)$
with samples from the latent prior. A cycle-consistency term penalises discrepancies between an
observed window and its reconstruction,
\[
  \mathcal{L}_{\mathrm{cyc}}
  =\mathbb{E}_{X}\!\left[\left\lVert X-G(E(X))\right\rVert_2\right].
\]

\begin{figure}[tbp]
  \centering
  \includegraphics[width=\linewidth]{swat/thesis/swat-gan-pipeline.pdf}
  \caption{Multivariate GAN detection pipeline. (a) The four networks are trained together on
  earlier windows containing normal and attack observations. (b) The later 393,679 test timestamps are
  scored with the trained networks; the latent critic has no part in the score $a_t$ of
  Equation~\ref{eq:anomaly:gan-score}. The blue tag marks the input representation, the only
  factor varied between runs, and orange tags mark development steps that used attack ground truth. The
  window length is fixed at 60 samples.}
  \label{fig:anomaly:gan-pipeline}
\end{figure}

Training combines the two Wasserstein objectives with this reconstruction constraint. The
encoder provides a direct mapping from process observations to latent representations, while
the cycle-consistency term encourages those representations to retain information required to
reconstruct the input. The data critic supplies an additional measure of how closely an observed
sequence resembles the learned training distribution. The latent critic contributes to training
and is not included in the anomaly score at inference \cite{geiger2020tadgan}.

At inference, each test window is encoded and reconstructed using the fitted model.
Discrepancies between the observed and reconstructed sequences provide a reconstruction score
(Figure~\ref{fig:anomaly:gan-reconstruction}). The data critic supplies a second score for the
observed window. Before combination, the scores require a common orientation, such that larger
values indicate greater deviation from the reference behaviour. Their normalisation must also
prevent a product of two negative standardised values from being interpreted as strong anomalous
evidence. The reported score therefore uses non-negative, consistently oriented normalised
components:
\begin{equation}
  a_t=\widetilde{r}_t\,\widetilde{c}_t.
  \label{eq:anomaly:gan-score}
\end{equation}

\begin{figure}[htbp]
  \centering
  \includegraphics[width=\linewidth]{swat/report-originals/report-fig-4.15-reconstruction.png}
  \caption{Observed signals (blue) and their reconstructions (orange) for two dimensions. The
  repeated red trace is the aggregate reconstruction discrepancy across dimensions, shared by
  the two displayed signals. It precedes the combination of
  Equation~\ref{eq:anomaly:gan-score}.}
  \label{fig:anomaly:gan-reconstruction}
\end{figure}

Score transformations are fitted on development data and fixed during evaluation. Reconstruction
discrepancy and critic output jointly indicate deviation, without providing a calibrated attack
probability. Normal operational changes can also produce high reconstruction error.

For each timestamp, overlapping-window reconstructions are averaged per variable before
calculating the multivariate reconstruction discrepancy; critic outputs are averaged separately.
Only windows within the same partition contribute. Each timestamp receives one combined score
and prediction and is counted once in the confusion matrix.

An adaptive threshold converts scores into binary predictions and may merge adjacent anomalous
intervals. Its historical context is distinct from the networks' 60-second input window.
Because aggregation includes observations later than the timestamp being scored, the results
describe offline detection; online performance and real-time detection delay remain unevaluated.

\subsection{Input Representations and Settings}
\label{sec:anomaly:gan-inputs}

The representation comparison includes the complete set of 51 process variables, two principal
component analysis (PCA) configurations, an autoencoder representation and a selected feature
subset. The full representation retains the original measurements and actuator states. It
provides the broadest input coverage but also exposes the model to redundant relationships and
normal operational variation across the process.

PCA projects the measurements onto $k$ orthogonal components. Let
$\lambda_1\geq\cdots\geq\lambda_{51}\geq0$ denote the eigenvalues of the covariance matrix
estimated from training data. The retained variance fraction is
\[
  \rho_k=\frac{\sum_{j=1}^{k}\lambda_j}{\sum_{j=1}^{51}\lambda_j}.
\]
The study compares $k=5$ and $k=24$, with $\rho_{24}\approx0.99$. Retained variance does not
quantify attack-relevant information; low-variance directions may still distinguish attacks.

The autoencoder learns a nonlinear representation for reconstruction before GAN processing.
Feature selection combines correlation analysis with chi-squared selection, recursive feature
elimination, LASSO and Random Forest importance. Supervised selectors use training labels only.
Any non-negative transformation required for chi-squared selection is fitted separately from GAN scaling.
The aggregate results do not specify the selected variables or the autoencoder bottleneck dimension.

All preprocessing is fitted on training data and reused for validation and testing. GAN inputs
use a training-derived scale with reference range $[-1,1]$. Values outside the training range
remain unclipped to preserve deviations.

For feature vectors $\mathbf{y}_t\in\mathbb{R}^{d}$, each window is
\[
  X_t=[\mathbf{y}_t,\ldots,\mathbf{y}_{t+59}]^{\mathsf T}
  \in\mathbb{R}^{60\times d}.
\]
At the nominal sampling rate, this gives a 60-second context with a one-second stride.
An uninterrupted sequence of $N\geq60$ observations yields
$N-60+1$ windows; performance metrics count timestamps after aggregation.

The final configuration uses Adam with learning rate $5\times10^{-4}$, latent dimension 20,
batch size 64 and 50 epochs. All variants share the evaluation cohort and counting convention.
Attribution of performance differences to the input representation requires control of network
capacity and tuning effort.

The reference computing platform comprises an AMD EPYC 7551P processor with 32 cores, 125\,GB of
RAM and a TITAN X GPU with 12\,GB of memory. These specifications describe the reported execution
environment. They do not establish training duration, inference throughput or memory
consumption. Reducing the input dimension from 51 variables to 24 components constitutes a
52.94\% reduction in input dimensionality; computational savings would require direct
measurement.

\subsection{Detection Results}
\label{sec:anomaly:gan-results}

Table~\ref{tab:anomaly:gan-representations} presents the measured performance of
the five input configurations. Each row uses the same class totals: 51,453 anomalous timestamps
and 342,226 normal timestamps. Precision describes the proportion of predicted anomalous
timestamps that fall within recorded attack periods. Recall describes the proportion of actual
anomalous timestamps that are detected. $F_1$ is the harmonic mean of precision and recall, while the
false-positive rate is the proportion of normal timestamps incorrectly flagged. All table entries
are calculated from integer confusion counts and rounded to two decimal places.

\begin{table}[tbp]
  \centering
  \caption{Measured timestamp-level performance of the input representations. Each timestamp
  is counted once and no point adjustment is applied.}
  \label{tab:anomaly:gan-representations}
  \small
  \begin{tabular}{@{}lrrrr@{}}
    \toprule
    Input representation & Precision (\%) & Recall (\%) & $F_1$ (\%) & False-positive rate (\%) \\
    \midrule
    All 51 variables           & 70.00 & 94.50 & 80.43 & 6.09 \\
    PCA: 5 components          & 75.00 & 86.50 & 80.34 & 4.34 \\
    PCA: 24 components         & 73.00 & 92.00 & 81.41 & 5.12 \\
    Autoencoder representation & 72.00 & 86.50 & 78.59 & 5.06 \\
    Selected feature subset    & 71.00 & 84.00 & 76.96 & 5.16 \\
    \bottomrule
  \end{tabular}
\end{table}

The complete representation has the highest recall, at 94.50\%, together with the
highest false-positive rate, at 6.09\%. This profile represents a detector that identifies a
larger proportion of anomalous observations while also flagging more normal observations. The
five-component PCA configuration has the highest precision, at 75.00\%, and the lowest
false-positive rate, at 4.34\%, accompanied by lower recall of 86.50\%. In these results,
stronger compression therefore corresponds to fewer false-positive timestamps and more
missed anomalous timestamps.

The 24-component PCA configuration provides an intermediate balance, with precision of 73.00\%,
recall of 92.00\% and an $F_1$ score of 81.41\%. Relative to the complete representation, its
precision is 3.00 percentage points higher and its recall is 2.50 percentage points lower. Its
$F_1$ score is approximately 0.98 percentage points higher. This configuration has the highest
reported $F_1$ and is used for the detailed analysis below. The numerical difference does not
establish a statistically reliable advantage or a preference under every operational cost model.

The autoencoder representation has a precision of 72.00\% and recall of 86.50\%,
corresponding to an $F_1$ score of 78.59\%. The selected feature subset has precision of
71.00\%, recall of 84.00\% and $F_1$ of 76.96\%. These values show lower detection coverage
for the two configurations. They do not identify a causal mechanism: loss of
relevant input information, differences in representation learning and changes in model
capacity would require separate investigation.

The final 24-component PCA evaluation produced the confusion matrix in
Table~\ref{tab:anomaly:gan-confusion}: 47,337 of the 51,453 anomalous timestamps are detected
and 17,508 of the 342,226 normal timestamps are flagged as anomalous.

\begin{table}[tbp]
  \centering
  \caption{Timestamp-level confusion matrix for the final 24-component PCA configuration. Counts
  refer to individual timestamps after aggregation of overlapping window outputs.}
  \label{tab:anomaly:gan-confusion}
  \small
  \begin{tabular}{@{}lrrr@{}}
    \toprule
    Actual class & Predicted anomalous & Predicted normal & Total \\
    \midrule
    Anomalous & 47,337 (TP) &   4,116 (FN) &  51,453 \\
    Normal    & 17,508 (FP) & 324,718 (TN) & 342,226 \\
    \midrule
    Total     & 64,845      & 328,834      & \textbf{393,679} \\
    \bottomrule
  \end{tabular}
\end{table}

The performance measures follow directly from these counts and are calculated from the
unrounded values:
\begin{alignat*}{3}
  \mathrm{Precision} &= TP/(TP+FP)          &&= 47{,}337/64{,}845   &&= 73.00\%,\\
  \mathrm{Recall}    &= TP/(TP+FN)          &&= 47{,}337/51{,}453   &&= 92.00\%,\\
  F_1                &= 2TP/(2TP+FP+FN)     &&= 94{,}674/116{,}298  &&= 81.41\%,\\
  \mathrm{FPR}       &= FP/(FP+TN)          &&= 17{,}508/342{,}226  &&= 5.12\%,\\
  \mathrm{Accuracy}  &= (TP+TN)/(TP+FN+FP+TN) &&= 372{,}055/393{,}679 &&= 94.51\%.
\end{alignat*}

The recall of 92.00\% means that the detector flags approximately 92 of every 100 anomalous
timestamps. It does not mean that 92\% of attack episodes are detected. Long attacks contribute
more timestamps than short attacks, so the same timestamp-level recall could arise from
different distributions of detection across attack episodes. Precision of 73.00\% means that
approximately 27\% of flagged timestamps fall within recorded normal periods. This proportion
differs from the false-positive rate of 5.12\%, whose denominator includes all normal
timestamps.

Accuracy is interpreted alongside the class distribution. Since normal observations account for
86.93\% of the cohort, a rule assigning every observation to the normal class would achieve
86.93\% accuracy while detecting no anomalous timestamps. The measured accuracy of 94.51\%
therefore provides supplementary information, whereas recall, precision and $F_1$ directly
describe the balance relevant to identifying the anomalous class.

No point adjustment is applied. A timestamp inside an attack interval counts as a true
positive only if that timestamp is predicted anomalous. Detecting one observation does not
automatically convert all other observations in the interval into true positives. Contiguous
positive predictions may be grouped into alarm intervals for display or operational analysis,
but grouping does not retrospectively alter the timestamp-level classifications used in the
matrix.

\subsection{Interpretation and Evaluation Scope}
\label{sec:anomaly:gan-scope}

The measured results describe a detector with high coverage of anomalous timestamps and a
substantial false-positive burden. At the nominal one-second cadence, 17,508 false-positive
observations correspond to approximately 4.86 hours of normal observations being
flagged. This duration does not identify the number of alarms presented to an operator. A
sustained interval and many isolated detections could contribute the same number of
false-positive timestamps while requiring very different responses. Likewise, the 4,116 missed
anomalous timestamps correspond to 68.6 minutes of anomalous observations, but their
operational significance depends on which attacks and phases of process disruption they
represent.

Attack-level assessment requires matching detected intervals to the retained attack episodes.
Relevant outcomes include the number of episodes detected, the time between attack onset and
the first qualifying detection, and the duration of detection within each episode. A minimum
overlap or persistence rule must be specified before those outcomes are calculated. Attacks
without a qualifying detection remain missed episodes and cannot be omitted from the
interpretation of detection delay. The denominator is the number of attack episodes present in
the evaluation cohort, which cannot be inferred solely from the 36 attacks described for the
source dataset.

False alarm episodes similarly require an explicit grouping rule and a normal-operation exposure
denominator. Their rate can be expressed as the number of distinct false alarm episodes per day
of evaluated normal operation. Neither this rate nor the distribution of alarm durations follows
from the confusion matrix alone. Timestamp-level precision also depends on the prevalence of
anomalous observations: its value in a cohort containing 13.07\% anomalous timestamps does not
establish the precision expected during deployment with a different attack prevalence.

The comparison between representations specifies a modest difference in $F_1$ rather than
evidence of a universal advantage for PCA. The complete input configuration retains the highest
recall, while the five-component PCA representation produces the fewest false-positive
timestamps relative to normal observations. The preferred configuration therefore depends on the
relative consequences of missed anomalous behaviour and unnecessary alarms. Retention of 99\% of
training variance provides a description of the 24-component representation, not a guarantee of
detection coverage or robustness.

The report gives no run-to-run variability from independently initialised training runs.
A repeated-run analysis under the same partition and model-selection procedure should calculate
precision, recall, $F_1$ and false-positive rate for each run, followed by their means and
standard deviations. Resampling-based uncertainty estimates must account for temporal dependence
and clustering within attack episodes because individual timestamps are not independent.

Further evaluation would separate the contributions of reconstruction scoring, critic scoring
and their combination, and examine window length under a common validation procedure.
Precision-recall curves would characterise changes across thresholds; their area cannot be
derived from one confusion matrix. Training duration, inference throughput and peak memory use
would establish computational cost on the stated hardware. The model topology, loss weights,
normalisation rule, threshold parameters, software versions, feature identities and partition
records would form the implementation record needed for reproducibility.

\section{Limitations, Validity, and Individual Contribution}
\label{sec:anomaly:limitations}

\paragraph{Scope of the industrial evaluation.}
The S7-300 result rests on 25 attack episodes recorded in one plant over 24 hours, and the
held-out partition holds a single episode of each class. Records taken from one episode are
correlated, so the per-class scores of Table~\ref{tab:anomaly:iuno-confusion} summarise five
attack events and not 1,233 independent trials. The benign monitoring day was a separate
exercise with four active PLCs and executed no attack, so it bounds the false-alarm rate and says
nothing about recall.
No experiment reported here trains on one cell and tests for attacks on another, so transfer
between installations is untested. The importances in
Figure~\ref{fig:anomaly:iuno-importance} are impurity decreases on the training partition, so
they describe what this forest used and do not show which fields detection requires. Both exercises were
labelled by the team that executed the attacks, and each two-hour baseline covers only the
operating modes that occurred within it.

\paragraph{Sampling and selection.}
The GAN study develops the model on the first 56,240 attack-period timestamps and evaluates it
on the later 393,679 timestamps. Training and validation also use separate normal-operation
data, and attack episodes do not cross partition boundaries. The evaluation remains within one
SWaT campaign. The S7-300 evaluation holds out whole capture sessions. Neither study tests transfer
to another installation or an operational class ratio. No repeated runs or uncertainty
intervals are reported, so small differences between configurations cannot be interpreted as
stable improvements.

\paragraph{Simple baselines.}
On SWaT, simple tests on process values have matched much more elaborate detectors under a
common protocol \cite{wolsing2022simple}. Reimplementing published intrusion detectors in one
framework has likewise left simple baselines competitive \cite{bilot2025sometimes}. No such
baseline is reported for the GAN study or for the S7-300 evaluation, so neither result shows
that its model outperforms a simple rule on the same data.

\paragraph{Individual contribution.}
The candidate designed and implemented \nadics{}, its model interface, and the S7comm and PROFINET
extractor. He produced the labelled S7-300 attack corpus. For the process study, he designed and
implemented the multivariate GAN pipeline, its preprocessing variants, dynamic architecture,
training, scoring, and evaluation on SWaT. He analysed the results and wrote the internal
report. His contribution to the exploratory service-graph study was the
concept and the measurement design. The IUNO deliverables attribute authorship at document level
\cite{iuno_ap4,iuno_ap7}, so this account rests on the candidate's attestation and on the
surviving artefacts rather than on passage-level attribution.

\section{Chapter Summary}
\label{sec:anomaly:summary}

RQ3 asks how a modular framework can analyse network and process features for anomaly
detection in industrial systems. \nadics{} provides a researcher-guided workflow for adapting
data preparation and comparing candidate learners through fixed interfaces. Its modular design
supports the selection of a configuration for a given dataset. On the S7-300 corpus recorded
inside IUNO the framework reaches 94.50
per cent attack $F_1$ and 89.56 per cent macro-$F_1$ over six classes at a record false-positive
rate of 0.484 per cent, and a benign day with four active PLCs raised nine alarm episodes, all of
them false. For SWaT process windows, the GAN evaluation reports
92.00 per cent timestamp recall at a false-positive rate of 5.12 per cent with 24 PCA
components. The full input gives the highest recall and five components the lowest
false-positive rate, so the preferred representation depends on the relative cost of missed
anomalies and false alarms.

The framework and its two industrial applications address the representation and modular-analysis
aspects of gap~\gapref{3}. Their results remain specific to the evaluated installation or SWaT
campaign; transfer to other installations and robustness against adaptive attacks remain
untested. Chapters~\ref{chap:advml} and~\ref{chap:discussion} examine requirements for broader
evaluation.
